\documentclass[10pt,a4paper]{amsart}
\usepackage{amsmath,amssymb,amsthm,graphicx}
\usepackage{verbatim}
\newenvironment{reference}{\comment}{\endcomment}
%\newenvironment{reference}{}{}
\newenvironment{slogan}{\comment}{\endcomment}
\newenvironment{history}{\comment}{\endcomment}
% For commutative diagrams we use Xy-pic
\usepackage[all]{xy}

% We use 2cell for 2-commutative diagrams.
\xyoption{2cell}
% Theorem environments.
%
\theoremstyle{plain}
\newtheorem{theorem}[subsection]{Theorem}
\newtheorem{proposition}[subsection]{Proposition}
\newtheorem{lemma}[subsection]{Lemma}

\theoremstyle{definition}
\newtheorem{definition}[subsection]{Definition}
\newtheorem{example}[subsection]{Example}
\newtheorem{exercise}[subsection]{Exercise}
\newtheorem{situation}[subsection]{Situation}

\theoremstyle{remark}
\newtheorem{remark}[subsection]{Remark}
\newtheorem{remarks}[subsection]{Remarks}

\numberwithin{equation}{subsection}

% Macros
%
\def\lim{\mathop{\mathrm{lim}}\nolimits}
\def\colim{\mathop{\mathrm{colim}}\nolimits}
\def\Spec{\mathop{\mathrm{Spec}}}
\def\Hom{\mathop{\mathrm{Hom}}\nolimits}
\def\Ext{\mathop{\mathrm{Ext}}\nolimits}
\def\SheafHom{\mathop{\mathcal{H}\!\mathit{om}}\nolimits}
\def\SheafExt{\mathop{\mathcal{E}\!\mathit{xt}}\nolimits}
\def\Sch{\mathit{Sch}}
\def\Mor{\mathop{\mathrm{Mor}}\nolimits}
\def\Ob{\mathop{\mathrm{Ob}}\nolimits}
\def\Sh{\mathop{\mathit{Sh}}\nolimits}
\def\NL{\mathop{N\!L}\nolimits}
\def\CH{\mathop{\mathrm{CH}}\nolimits}
\def\proetale{{pro\text{-}\acute{e}tale}}
\def\etale{{\acute{e}tale}}
\def\QCoh{\mathit{QCoh}}
\def\Ker{\mathop{\mathrm{Ker}}}
\def\Im{\mathop{\mathrm{Im}}}
\def\Coker{\mathop{\mathrm{Coker}}}
\def\Coim{\mathop{\mathrm{Coim}}}

% Boxtimes
%
\DeclareMathSymbol{\boxtimes}{\mathbin}{AMSa}{"02}

%
% Macros for moduli stacks/spaces
%
\def\QCohstack{\mathcal{QC}\!\mathit{oh}}
\def\Cohstack{\mathcal{C}\!\mathit{oh}}
\def\Spacesstack{\mathcal{S}\!\mathit{paces}}
\def\Quotfunctor{\mathrm{Quot}}
\def\Hilbfunctor{\mathrm{Hilb}}
\def\Curvesstack{\mathcal{C}\!\mathit{urves}}
\def\Polarizedstack{\mathcal{P}\!\mathit{olarized}}
\def\Complexesstack{\mathcal{C}\!\mathit{omplexes}}
% \Pic is the operator that assigns to X its picard group, usage \Pic(X)
% \Picardstack_{X/B} denotes the Picard stack of X over B
% \Picardfunctor_{X/B} denotes the Picard functor of X over B
\def\Pic{\mathop{\mathrm{Pic}}\nolimits}
\def\Picardstack{\mathcal{P}\!\mathit{ic}}
\def\Picardfunctor{\mathrm{Pic}}
\def\Deformationcategory{\mathcal{D}\!\mathit{ef}}
\usepackage{cleveref}
\crefname{originalreference}{original reference}{original references}
\Crefname{originalreference}{Original reference}{Original references}
\makeatletter
\newlabel{homology-lemma-check-exactness}{{homology-lemma-check-exactness}{source}{}{sourcepointer-1}{}}
\newlabel{homology-lemma-check-exactness@cref}{{[originalreference][0][]homology-lemma-check-exactness}{[1][0][]source}{}{}{}}
\makeatother
\title{The sheaf of homomorphisms is exact in its first variable on the right and in its second variable on the left}
\author{The Stacks Project Authors}
\date{stacks-project:modules.tex:local-snapshot-20261012; GNU FDL 1.2 or later}
\begin{document}
\maketitle
\noindent\textit{Editorial note.} All delivered mathematical text is the author's own, copied from the pinned Stacks Project chapter source sources/stacks-modules/pinned-original-source.tex; no statement or proof step was rewritten, reordered or removed.\par\medskip
\noindent\textit{Editorial note.} Preamble provenance: the Stacks Project keeps one single shared preamble for every chapter; this card takes it directly from the staged sources/stacks-modules/preamble.tex as three verbatim spans (lines 8-12, 14-18 and 154-217), so no derived file was written and nothing was added or redefined. The wrapper class and the base packages are supplied by the assembler.\par\medskip
\noindent\textit{Editorial note.} Bibliography provenance: neither the counted statement nor its proof contains a citation (this lane rescans the whole span before writing the card), so the card carries no bibliography and no entry was reconstructed.\par\medskip
\noindent\textit{Editorial note.} Selection provenance: this card was opened under the file-based control contract (CONTROL.md), with the mechanical pre-card gate run on its statement and proof interval before the card was written.\par\medskip
\noindent\textit{Allowed standard prerequisites (editorial).}\begin{itemize}
\item The sheaf of homomorphisms between two modules and the exactness of its sections over an open as a sequence of modules. The author's section on the internal hom, in which the counted statement occurs. Fixes the sheaves whose exactness in the two displays is to be proved.
\item The criterion that a sequence of sheaves is exact if and only if its sections over every open form an exact sequence. Invoked inside the counted proof at source lines 3808-3809 as a cross-chapter pointer. Converts the exactness of the sequences of sections into exactness of the displayed sequences of sheaves.
\end{itemize}\medskip
\ifdefined\hypertarget\hypertarget{sourcepointer-1}{}\fi
\noindent\textit{Source reference.} Original reference homology-lemma-check-exactness: sources/stacks-modules/pinned-original-source.tex, source lines 3809--3809. Navigation pointer only: the counted text refers to the Stacks item homology-lemma-check-exactness, which belongs to another chapter of the project and is not part of this card. The pointer carries the target's own label as its number because this chapter's staged snapshot does not contain the other chapter's tag table; its content is not reproduced here.\par\medskip
\begin{lemma}
\label{lemma-internal-hom-exact}
Let $(X, \mathcal{O}_X)$ be a ringed space.
Let $\mathcal{F}$, $\mathcal{G}$ be $\mathcal{O}_X$-modules.
\begin{enumerate}
\item If $\mathcal{F}_2 \to \mathcal{F}_1 \to \mathcal{F} \to 0$
is an exact sequence of $\mathcal{O}_X$-modules, then
$$
0 \to
\SheafHom_{\mathcal{O}_X}(\mathcal{F}, \mathcal{G}) \to
\SheafHom_{\mathcal{O}_X}(\mathcal{F}_1, \mathcal{G}) \to
\SheafHom_{\mathcal{O}_X}(\mathcal{F}_2, \mathcal{G})
$$
is exact.
\item If $0 \to \mathcal{G} \to \mathcal{G}_1 \to \mathcal{G}_2$
is an exact sequence of $\mathcal{O}_X$-modules, then
$$
0 \to
\SheafHom_{\mathcal{O}_X}(\mathcal{F}, \mathcal{G}) \to
\SheafHom_{\mathcal{O}_X}(\mathcal{F}, \mathcal{G}_1) \to
\SheafHom_{\mathcal{O}_X}(\mathcal{F}, \mathcal{G}_2)
$$
is exact.
\end{enumerate}
\end{lemma}

\begin{proof}
Let $\mathcal{F}_2 \to \mathcal{F}_1 \to \mathcal{F} \to 0$
be as in (1). For every $U \subset X$ open the sequence
$$
0 \to \Hom_{\mathcal{O}_U}(\mathcal{F}|_U, \mathcal{G}|_U) \to
\Hom_{\mathcal{O}_U}(\mathcal{F}_1|_U, \mathcal{G}|_U) \to
\Hom_{\mathcal{O}_U}(\mathcal{F}_2|_U, \mathcal{G}|_U)
$$
is exact by Homology, Lemma \ref{homology-lemma-check-exactness}.
This means that
taking sections over $U$ of the sequence of sheaves in (1) produces
an exact sequence of abelian groups. Hence the sequence in (1) is exact
by definition. The proof of (2) is exactly the same.
\end{proof}

\end{document}
